\subsection{Simple compactness of sets of continuous functions}

In this section, X denotes a compact space and \(\mathcal{C}_{s}(X)\) the space of continuous functions on X, with values in the field K (equal to R or C). The space \(\mathcal{C}_{s}(X)\) is assigned the topology of simple convergence on X.

PROPOSITION 2. --- Let D be a dense subset of X and A a subset of the space \(\mathcal{C}_{s}(X)\) . The following conditions are equivalent :

\begin{enumerate}
\def\labelenumi{(\roman{enumi})}
\item
  A is relatively compact in \(\mathcal{C}_s(\mathbf{X})\) .
\item
  From every infinite sequence of elements of \(\mathbf{A}\) , we can extract a sequence converging in \(\mathcal{C}_s(\mathbf{X})\) .
\item
  Every infinite sequence of elements of \(\mathbf{A}\) has a limit point in \(\mathcal{C}_s(\mathbf{X})\) .
\item
  Let \((f_n)_{n\in \mathbb{N}}\) be a sequence of functions belonging to A and \((x_m)_{m\in \mathbb{N}}\) a sequence of points of D. If the iterated limits
\end{enumerate}

\[
\gamma = \lim _ {m \rightarrow \infty} \lim _ {n \rightarrow \infty} f _ {n} (x _ {m}), \quad \delta = \lim _ {n \rightarrow \infty} \lim _ {m \rightarrow \infty} f _ {n} (x _ {m})
\]

exist, then they are equal. In addition, we have \(\sup_{f\in \mathbf{A}}|f(x)| < +\infty\) for all \(x\in \mathbf{X}\) .

\begin{enumerate}
\def\labelenumi{(\roman{enumi})}
\tightlist
\item
  \(\Rightarrow\) (ii): let \(\overline{\mathbf{A}}\) be the closure of \(\mathbf{A}\) in \(\mathcal{C}_s(\mathbf{X})\) . Assume that \(\mathbf{A}\) is compact, and consider a sequence of functions \(f_n \in \mathbf{A}\) (for \(n \in \mathbf{N}\) ). Let \(\phi\) be the continuous mapping \(x \mapsto (f_n(x))_{n \in \mathbf{N}}\) from \(\mathbf{X}\) into the metrizable space \(\mathbf{K}^{\mathbf{N}}\) . The image \(\mathbf{X}'\) of \(\mathbf{X}\) under \(\phi\) is a compact metrizable space, since \(\mathbf{X}\) is compact. Let \(\mathbf{E}\) be the closed subspace of \(\mathcal{C}_s(\mathbf{X})\) consisting of continuous functions \(f\) on \(\mathbf{X}\) such that the relation \(\phi(x) = \phi(y)\) implies \(f(x) = f(y)\) for every pair of points \(x, y\) in \(\mathbf{X}\) . By cor. 2 of GT, I, § 9, No.~4 and prop. 3 of GT, I, § 5, No.~2, the mapping \(f' \mapsto f' \circ \phi\) is a homeomorphism \(\phi^*\) from \(\mathcal{C}_s(\mathbf{X}')\) onto E. Hence the set \(\mathbf{A}' = (\phi^*)^{-1}(\overline{\mathbf{A}})\) is compact in \(\mathcal{C}_s(\mathbf{X}')\) , and it is clear that there exist elements \(f_n'\) in \(\mathbf{A}'\) such that \(\phi^*(f_n') = f_n' \circ \phi\) is equal to \(f_n\) .
\end{enumerate}

Since \(X'\) is a compact metrizable space, there exists a countable dense subset \(D'\) in \(X'\) (GT, IX, §2, No.~8, prop. 12 and §2, No.~9, prop. 16). Let \(T_{1}\) (resp. \(T_{2}\) ) be the topology on \(A'\) induced by the topology of simple convergence on \(D'\) (resp. \(X'\) ). Then \(T_{1}\) is metrizable, \(T_{2}\) is compact and finer than \(T_{1}\) , hence \(T_{1}\) and \(T_{2}\) coincide; in other words, \(A'\) is a compact metrizable subspace of \(\mathcal{C}_{s}(X')\) . Therefore, there exists a sequence \((f_{n_k}^{\prime})\) extracted from \((f_n^{\prime})\) and converging to an element \(f^{\prime}\) of \(\mathcal{C}_s(\mathbf{X}')\) . Therefore, the sequence \((f_{n_k})\) converges to \(f = f' \circ \phi\) in \(\mathcal{C}_s(\mathbf{X})\) .

\begin{enumerate}
\def\labelenumi{(\roman{enumi})}
\setcounter{enumi}{1}
\item
  \(\Rightarrow\) (iii): this is clear.
\item
  \(\Rightarrow\) (iv): suppose that every infinite sequence of elements of A has a limit point in \(\mathcal{C}_s(\mathbf{X})\) . Let \(x \in \mathbf{X}\) . The mapping \(\phi_x: f \mapsto f(x)\) from A into K is continuous. Consequently, every infinite sequence in \(\phi_x(\mathbf{A})\) has a limit point; since the field K (equal to R or C) is metrizable, the set \(\phi_x(\mathbf{A})\) is relatively compact in K, hence bounded. In other words, we have \(\sup_{f \in \mathbf{A}} |f(x)| < \infty\) .
\end{enumerate}

Let \(f_n, x_m, \gamma\) and \(\delta\) be as in (iv). Let \(f\) be a limit point of the sequence \((f_n)\) in \(\mathcal{C}_s(\mathbf{X})\) , and let \(x\) be a limit point of the sequence \((x_m)\) in the compact space \(\mathbf{X}\) . For every \(m\) , the mapping \(h \mapsto h(x_m)\) from \(\mathcal{C}_s(\mathbf{X})\) into \(\mathbf{K}\) is continuous. In view of the hypotheses, we have \(f(x_m) = \lim_{n \to \infty} f_n(x_m)\) , and hence \(\gamma = \lim_{m \to \infty} f(x_m)\) ; since \(f: \mathbf{X} \to \mathbf{K}\) is continuous and \(x\) is a limit point of the sequence \((x_m)\) , we get \(\gamma = f(x)\) . In an analogous way, we prove the equality \(\delta = f(x)\) , whence \(\gamma = \delta\) .

\begin{enumerate}
\def\labelenumi{(\roman{enumi})}
\setcounter{enumi}{3}
\tightlist
\item
  \(\Rightarrow\) (i): suppose that the set of numbers \(f(x)\) , as \(f\) ranges over \(\mathbf{A}\) , is bounded in \(\mathbf{K}\) for all \(x\in \mathbf{X}\) . This is equivalent to assuming that the closure \(\overline{\mathbf{A}}\) of \(\mathbf{A}\) in the product space \(\mathbf{K}^{\mathbf{X}}\) is compact (GT, I, § 9, No.~5). Suppose that \(\mathbf{A}\) is not relatively compact in \(\mathcal{C}_s(\mathbf{X})\) . This means that there exists a function \(u\in \overline{\mathbf{A}}\) and a point \(a\in \mathbf{X}\) such that \(u\) is not continuous at \(a\) . Hence there exists a real number \(\varepsilon >0\) such that in every neighbourhood \(\mathrm{U}\) of \(a\) , there exists a point \(x\) with \(|u(x) - u(a)|\geqslant \varepsilon\) .
\end{enumerate}

We shall construct by induction a sequence \((x_{n})_{n\in \mathbf{N}}\) of points in \(\mathbf{D}\) and a sequence \((f_n)_{n\in \mathbf{N}}\) of elements of \(\mathbf{A}\) , satisfying the following relations:

\[
\left| u (x _ {m}) - u (a) \right| \geqslant \varepsilon \quad \text { for } \quad m \geqslant 1;\tag{$(1)_{m}$}
\]

\[
\left| u (x _ {i}) - f _ {m} (x _ {i}) \right| \leqslant \frac {1}{m + 1} \text { for } 0 \leqslant i \leqslant m - 1;\tag{$(2)_{m}$}
\]

\[
\left| f _ {m} (x _ {i}) - f _ {m} (a) \right| \leqslant \frac {1}{i + 1} \quad \text { for } \quad 0 \leqslant m \leqslant i.\tag{$(3)_{m,i}$}
\]

We take \(x_0 = a\) with \(f_0\) arbitrary in A (the set A is not empty, otherwise it will be relatively compact in \(\mathcal{C}_s(\mathbf{X})\) ). Let \(n \geqslant 1\) and \(x_0, x_1, \ldots, x_{n-1}, f_0, f_1, \ldots, f_{n-1}\) satisfy relations \((1)_m, (2)_m\) for \(1 \leqslant m < n\) and \((3)_{m,i}\) for \(0 \leqslant m \leqslant i < n\) . Since \(u\) belongs to \(\overline{\mathbf{A}}\) , there exists \(f_n \in \mathbf{A}\) satisfying \((2)_n\) . Let \(V_n\) be the set of all \(x \in X\) such that \(|f_m(x) - f_m(a)| \leqslant \frac{1}{n+1}\) for \(0 \leqslant m \leqslant n\) . Since \(f_n\) is continuous, \(V_n\) is a neighbourhood of \(a\) ; choose a point \(x_n\) in \(D \cap V_n\) such that \(|u(x_n) - u(a)| \geqslant \varepsilon\) , hence \((1)_n\) and \((3)_{m,n}\) are satisfied. Therefore, the construction can be continued.

Since \(u(\mathbf{X})\) is a compact subset of \(\mathbf{K}\) , there exists a sequence \((y_k)\) extracted from \((x_m)\) and such that the limit \(\gamma = \lim_{k\to \infty}u(y_k)\) exists. By (2) \(_m\) , we have \(u(x_i) = \lim_{n\to \infty}f_n(x_i)\) for all \(i\in \mathbf{N}\) , hence

\[
\gamma = \lim _ {k \rightarrow \infty} \lim _ {n \rightarrow \infty} f _ {n} (y _ {k}).\tag{4}
\]

On the other hand we have \(f_{n}(a) = \lim_{i\to \infty}f_{n}(x_{i})\) by (3) \(_{m,i}\) hence \(f_{n}(a) = \lim_{k\to \infty}f_{n}(y_{k})\) . Since \(x_0 = a\) , we deduce from (2) \(_{m}\) that \(\lim_{n\to \infty}f_n(a) = u(a)\) . Consequently,

\[
u (a) = \lim _ {n \rightarrow \infty} \lim _ {k \rightarrow \infty} f _ {n} (y _ {k}).\tag{5}
\]

Finally, from \((1)_{m}\) , we get \(|\gamma - u(a)| \geqslant \varepsilon\) , and so \(\gamma \neq u(a)\) . This contradicts assertion (iv); we have thus proved that (iv) implies (i).

